% TikZ figure for tool overview
\begin{tikzpicture}[
    node distance=0.95cm and 0.55cm,
    box/.style={rectangle, draw, line width=0.8pt, minimum width=1.45cm, minimum height=0.95cm, align=center, font=\Huge},
    widebox/.style={rectangle, draw, line width=0.8pt, minimum width=1.9cm, minimum height=0.95cm, align=center, font=\Huge},
    arrow/.style={->, line width=0.8pt, >=stealth}
]

% Define main flow nodes (centered vertically)
\node[box, fill=blue!20] (input) {Input Specs \\ (Base, Size)};
\node[box, fill=green!20, right=of input] (tiler) {Tiler};
\node[box, fill=yellow!20, right=of tiler] (combgen) {Decomposition\\ Scheme \\Combinations};

% Exporter and Visualizer stacked (centered around the main flow)
\node[box, fill=orange!20, right=of combgen, xshift=0.25cm, yshift=0.8cm] (exporter) {Intermediate\\Represetnation};
\node[box, fill=purple!20, right=of combgen, xshift=0.25cm, yshift=-1.5cm] (visualizer) {Visualizer};

% Visualization icon output from Visualizer (tiling grid representation)
\node[draw, rectangle, line width=0.7pt, inner sep=2pt, minimum height=1.45cm, minimum width=1.45cm, right=of visualizer, xshift=0.25cm] (pdfout) {
    \begin{tikzpicture}[scale=0.5]
        % Draw a simple grid pattern to represent tiling
        \draw[line width=0.4pt] (0,0) grid[step=0.5] (2,2);
        % Add some colored blocks to represent different tile types
        \fill[blue!30] (0,1) rectangle (0.5,1.5);
        \fill[green!30] (0.5,1) rectangle (1,1.5);
        \fill[red!30] (1,1) rectangle (1.5,1.5);
        \fill[yellow!30] (1.5,1) rectangle (2,1.5);
    \end{tikzpicture}
};

% Tester and Compiler stacked (centered around exporter vertically)
\node[widebox, fill=red!20, right=of exporter, xshift=-0.35cm, yshift=0.8cm] (tester) {Software\\Simulation};

% Verilog Generator
\node[box, fill=lime!20, right=of exporter, xshift=-1.40cm, yshift=-1.6cm] (verilog) {Verilog\\Generator};

% Hardware Simulator (to the right of Verilog Generator)
\node[widebox, fill=teal!20, right=of verilog, xshift=0.15cm] (hwsim) {Hardware\\Simulation};

% Output file icons (below Verilog Generator)
\node[draw, rounded corners=1pt, fill=white, minimum width=1.35cm, minimum height=0.72cm, below=1.0cm of verilog, xshift=-1.45cm, font=\Large] (file1) {.v};
\node[draw, rounded corners=1pt, fill=white, minimum width=1.35cm, minimum height=0.72cm, below=1.0cm of verilog, font=\Large] (file2) {.tcl};
\node[draw, rounded corners=1pt, fill=white, minimum width=1.35cm, minimum height=0.72cm, below=1.0cm of verilog, xshift=1.45cm, font=\Large] (file3) {.xdc};

% Test result icons (force horizontal alignment with software-simulation output)
\coordinate (testcenter) at (hwsim.center |- tester.center);
\node[draw, circle, fill=green!40, minimum size=1.05cm, at={($(testcenter)+(-0.72cm,0)$)}, font=\Large] (testpass) {\checkmark};
\node[draw, circle, fill=red!40, minimum size=1.05cm, at={($(testcenter)+(0.72cm,0)$)}, font=\Large] (testfail) {$\times$};

% Bounding box around test results
\node[draw, rectangle, line width=0.8pt, fit={(testpass) (testfail)}, inner sep=0.22cm] (testbox) {};

% Define arrows
\draw[arrow] (input) -- (tiler);
\draw[arrow] (tiler) -- (combgen);

% Arrow from combgen that splits to exporter and visualizer
\draw[thick] (combgen.east) -- ++(0.55,0) coordinate (split1);
\draw[arrow] (split1) |- (exporter.west);
\draw[arrow] (split1) |- (visualizer.west);

% Arrow from visualizer to PDF output
\draw[arrow] (visualizer.east) -- (pdfout.west);

% Arrow from exporter that splits to tester and verilog
\draw[thick] (exporter.east) -- ++(0.55,0) coordinate (split2);
\draw[arrow] (split2) |- (tester.west);
\draw[arrow] (split2) |- (verilog.west);

% Arrow from tester to test results bounding box
\draw[arrow] (tester.east) -- (testbox.west);

% Arrow from verilog generator to hardware simulator
\draw[arrow] (verilog) -- (hwsim);

% Arrow from hardware simulator to test results box (straight up)
\draw[arrow] (hwsim) -- (testbox);

% Arrows from verilog generator to output files
\draw[arrow] (verilog) -- (file1);
\draw[arrow] (verilog) -- (file2);
\draw[arrow] (verilog) -- (file3);

\end{tikzpicture}
